\documentclass[14pt, b4paper]{extreport}

\usepackage[
margin=1in
%top=1in,
%bottom=1in,
%left=0.5in,
%right=0.5in
]{geometry}

% For the title page
\usepackage{xcolor} 
\usepackage{graphicx}
\usepackage{ifthen}
\usepackage{tikz}
\usetikzlibrary{decorations.markings}
\usetikzlibrary{arrows.meta}
\usetikzlibrary{hobby, bending}
\tikzset{
	mydash/.style={dashed, dash pattern=on 6pt off 2pt},
	arrow inside/.style 2 args = {
		postaction={
			decorate,
			decoration={
				markings,
				mark=at position {#1*\pgfdecoratedpathlength-0.15cm} with {\coordinate (bent arrow 1);},
				mark=at position {#1*\pgfdecoratedpathlength-0.05cm} with {\coordinate (bent arrow 2);},
				mark=at position {#1*\pgfdecoratedpathlength+0.05cm} with {\coordinate (bent arrow 3);},
				mark=at position {#1*\pgfdecoratedpathlength+0.15cm} with {
					\coordinate (bent arrow 4);
					\ifthenelse{\equal{#2}{for}}
					{\draw[-{Triangle[length=3mm, width=1.6mm, bend]}] 
						(bent arrow 1) to[curve through={(bent arrow 2) .. (bent arrow 3)}] (bent arrow 4);
					}
					{
						\draw[-{Triangle[length=3mm, width=1.6mm, bend]}] 
						(bent arrow 4) to[curve through={(bent arrow 3) .. (bent arrow 2)}] (bent arrow 1);
					}
				}
			}
		}
	}
}

\def\ANS#1{\hspace*{2em}\textit{ج}: #1}
\def\SUG#1{\hspace*{2em}\textit{اقتراح}: #1}

%\usepackage{indentfirst}
\usepackage{float}

\newcounter{myfig}
\makeatletter
\newcommand{\myfig}[1][]{
	\refstepcounter{myfig}
	\textbf{الشكل \themyfig}
	\ifx\@empty#1\@empty
	\else
	\\[-2pt]#1
	\fi
}
\makeatother
\newcommand{\FIGURE}[1]{\noindent\hspace{20pt}\parbox[b]{3cm}{\myfig[#1]}\\}


% For changing the headers formatting
\usepackage{titletoc, titlesec}
\usepackage[inline]{enumitem}
\setlist{label=\textbf{\arabic*.}}
%\settasks{label-offset=25pt,label=(\xslalph*)}
\usepackage{environ}
\NewEnviron{tasks}
{
	\begin{enumerate*}[label=(\xslalph*), itemjoin=\hfil]
		\BODY
	\end{enumerate*}
}

\NewEnviron{tasks*}
{
	\begin{enumerate*}[
		before=\noindent,
		topsep=0pt,
		partopsep=0pt,
		parsep=5pt,
		itemsep=0pt,
		itemjoin=\\,
		label=(\xslalph*)\,\,,
		mode=unboxed
		]
		\BODY
	\end{enumerate*}
}

\makeatletter
\newcommand{\xslalph}[1]{\expandafter\@xslalph\csname c@#1\endcsname}
\newcommand{\@xslalph}[1]{%
	\ifcase#1
	\or أ%
	\or ب%
	\or جـ%
	\or د%
	\or هـ%
	\or و%
	\or ز%
	\or حـ%
	\or ط%
	\or ي%
	\or ك%
	\or ل%
	\or م%
	\or ن%
	\or س%
	\or ع%
	\or ف%
	\or ص%
	\or ق%
	\or ر%
	\or ش%
	\or ت%
	\or ث%
	\or خـ%
	\or ذ%
	\or ض%
	\or ظ%
	\or غ%
	\else \@ctrerr
	\fi
}
\AddEnumerateCounter{\xslalph}{\@xslalph}{m}
\makeatother

% For math
\usepackage{amsmath, amssymb, amsthm, needspace, amsfonts}

\newcommand{\conj}[1]{\mkern 0.5mu\overline{\mkern-0.5mu#1\mkern-0.5mu}\mkern 0.5mu}

\allowdisplaybreaks

% To use the arabic language
\usepackage{polyglossia}

\usepackage[]{setspace}
\setstretch{2}

\usepackage{tocbasic}
\addtotoclist[report.cls]{toc}
\renewcommand{\tableofcontents}{\listoftoc[{\contentsname}]{toc}}
\AfterTOCHead[toc]{\thispagestyle{empty}\pagestyle{empty}}
\AfterStartingTOC[toc]{\clearpage}


%\renewcommand{\thepage}{{\arabic{page}}}


% For fancy headers
\usepackage[twoside]{fancyhdr}
\pagestyle{fancy}
\fancyhf{}
\fancyhead[RO]{\ar{\leftmark\qquad\qquad\textbf{\thepage}}}
\fancyhead[RE]{\ar{\rightmark\qquad\qquad\textbf{\thepage}}}
\fancyhead[LO]{\small\ar{الفصل \thechapter}}
\fancyhead[LE]{\small{\ar{\textbf{بــنــد \arabic{section}}}}}

\renewcommand{\chaptermark}[1]{\markboth{#1}{}}
\renewcommand{\sectionmark}[1]{\markright{#1}{}}
\renewcommand{\headrulewidth}{0pt}
\setlength{\headsep}{3em}

\fancypagestyle{plain}{
	\fancyhf{}
	\fancyfoot[L]{\textbf{\thepage}}
}


% Set the arabic language as default and english as side langiage
\setmainlanguage[numerals=maghrib]{arabic}
\setotherlanguage{english}

% Set the arabic and the english font
%\newfontfamily{\amiritimes}{Amiri Times}[Script=Arabic]
\newfontfamily\titlefont{DecoType Naskh Swashes}[Script=Arabic]
\newfontfamily\arabicfont[Script=Arabic]{Amiri}
\newfontfamily{\farisifont}[Script=Arabic]{Noto Nastaliq Urdu}
\newfontfamily{\timesfont}{Times New Roman}
\rightfootnoterule
\renewcommand{\footnotesize}{\small}



\usepackage[T1]{fontenc}
\usepackage{times}

% Aliases for fast inserting arabic or english text
\newcommand{\ar}{\textarabic}
\newcommand{\en}{\textenglish}

% Load the MathTime Proffesional package
\usepackage[lite, subscriptcorrection, zswash]{mtpro2}
\straightbraces

\DeclareMathSymbol{0}{\mathalpha}{operators}{`0}
\DeclareMathSymbol{1}{\mathalpha}{operators}{`1}
\DeclareMathSymbol{2}{\mathalpha}{operators}{`2}
\DeclareMathSymbol{3}{\mathalpha}{operators}{`3}
\DeclareMathSymbol{4}{\mathalpha}{operators}{`4}
\DeclareMathSymbol{5}{\mathalpha}{operators}{`5}
\DeclareMathSymbol{6}{\mathalpha}{operators}{`6}
\DeclareMathSymbol{7}{\mathalpha}{operators}{`7}
\DeclareMathSymbol{8}{\mathalpha}{operators}{`8}
\DeclareMathSymbol{9}{\mathalpha}{operators}{`9}

%\makeatletter
%\renewcommand{\@biblabel}[1]{{$\selectlanguage[numerals=maghrib]{arabic} \mathbf{[#1]}$}}
%\renewcommand{\@makefntext}[1]{%
%	\noindent\makebox[2em][r]{\@thefnmark }#1
%}
%\makeatother

\newcommand{\arabicword}[1]{%
	\ifcase#1
	\or الأول
	\or الثاني
	\or الثالث
	\or الرابع
	\or الخامس 
	\or السادس
    \or السابع
	\or الثامن
	\or التاسع
	\or العاشر
	\else \number#1
	\fi
}

\renewcommand{\thechapter}{\arabicword{\value{chapter}}}

% Change the chapter header format
\makeatletter
\def\@makechapterhead#1{%
	\noindent
	\rule{\textwidth}{2pt}\\[10pt]
	{\huge \textbf{الفصل \thechapter}}\\[20pt]
	{\Huge \textbf{#1}}\\[10pt]
	\rule{\textwidth}{2pt}\\[9em]
%	\thispagestyle{empty}   % No page number
%	\newpage
}

\def\@makeschapterhead#1{
	\noindent
	\rule{\textwidth}{2pt}\\[15pt]
	{\huge \textbf{#1}}\\[10pt]
	\rule{\textwidth}{2pt}\\[5em]
}

\@removefromreset{section}{chapter}

%\def\tagform@#1{\maketag@@@{$(#1)$}}
\makeatother


% Chnage the section header format
\titleformat{\section}[block]
{\bfseries\fontsize{20}{17}\selectfont} 
{\thesection}{0.4em}{}


%\titlespacing{\section}{\parindent}{3em}{5pt}

% Change the subsection header format
\titleformat{\subsection}[block]
{\bfseries\fontsize{16}{17}\selectfont} 
{\thesubsection}{1em}{}  



\renewcommand{\thesection}{{\arabic{section}}.}

\renewcommand{\thesubsection}{{\arabic{chapter} - \arabic{section} - \arabic{subsection}}}

\titlecontents{chapter}[0em]{\large\bfseries\vspace{15pt}}{الفصل \thecontentslabel: }{}{\qquad\contentspage}

\titlecontents{chapter*}[0em]{\bfseries\vspace{15pt}}{}{}{\hfill\normalfont{\contentspage}}

\titlecontents{section}[6em]{\normalfont}{}{}{\qquad\contentspage}


\renewcommand{\thetable}{\arabic{chapter} - \arabic{table}}

\renewcommand{\thefigure}{{\textbf{\arabic{figure}}}}


\newtheoremstyle{mystyle}
{\topsep}
{\topsep}
{}
{}
{\bfseries}
{}
{\newline}  
{#1 #2 \textbf{#3}}

\theoremstyle{mystyle}

\newtheorem{theorem}{مبرهنة}[section]
\newtheorem{definition}{تعريف}[section]
\newtheorem{example}{مثال}[section]
\newtheorem{corollary}[theorem]{نتيجة}
%\newtheorem*{note}{ملاحظة}


\renewcommand{\thetheorem}{{\arabic{chapter} - \arabic{section} - \arabic{theorem}}}
\renewcommand{\thedefinition}{{\arabic{chapter} - \arabic{section} - \arabic{definition}}}
\renewcommand{\theexample}{{\arabic{chapter} - \arabic{section} - \arabic{example}}}
\renewcommand{\thecorollary}{{\arabic{chapter} - \arabic{section} - \arabic{corollary}}}

\newenvironment{myproof}{\noindent\textbf{البرهان}\vspace{5pt}\\ \noindent}{\hfill \(\qedsymbol\)}

\newenvironment{solution}{\noindent\textbf{الحل}\vspace{5pt}\\}{}
\newenvironment{note}{\noindent\textbf{ملاحظة}}{}
\newenvironment{notes}{\noindent\textbf{ملاحظات}\vspace{5pt}\\}{}

\pretocmd{\section}{\Needspace{5\baselineskip}}{}{}
\pretocmd{\subsection}{\Needspace{5\baselineskip}}{}{}
\AtBeginEnvironment{definition}{\Needspace{3\baselineskip}}
\AtBeginEnvironment{theorem}{\Needspace{3\baselineskip}}
\AtBeginEnvironment{example}{\Needspace{3\baselineskip}}
\AtBeginEnvironment{solution}{\Needspace{3\baselineskip}}
\AtBeginEnvironment{note}{\Needspace{3\baselineskip}}
\AtBeginEnvironment{myproof}{\Needspace{3\baselineskip}}

\newcommand{\N}{\mathbb{N}}
\newcommand{\Z}{\mathbb{Z}}
\newcommand{\Q}{\mathbb{Q}}
\newcommand{\R}{\mathbb{R}}
\newcommand{\C}{\mathbb{C}}
\newcommand{\LL}{\mathcal{L}}

% Define new math operators
\DeclareMathOperator{\Log}{Log}
\DeclareMathOperator{\lcm}{lcm}
\DeclareMathOperator{\arcsec}{arcsin}
\DeclareMathOperator{\arccsc}{arccsc}
\DeclareMathOperator{\arccot}{arccot}
\DeclareMathOperator{\re}{Re}
\DeclareMathOperator{\im}{Im}

\renewcommand{\mathbf}[1]{\textbf{\en{#1}}}


